\hypertarget{ux43eux431ux43eux441ux43dux43eux432ux430ux43dux438ux435-ux437ux43dux430ux447ux438ux43cux43eux441ux442ux438}{%
\chapter{Обоснование
значимости}\label{ux43eux431ux43eux441ux43dux43eux432ux430ux43dux438ux435-ux437ux43dux430ux447ux438ux43cux43eux441ux442ux438}}

Современные операционные системы являются основой функционирования любой
вычислительной техники --- от персональных компьютеров до серверных
кластеров и мобильных устройств. Понимание принципов и видов загрузки ОС
важно не только для системных администраторов, но и для всех
специалистов, связанных с IT-инфраструктурой. Изучение механизмов
BIOS/UEFI, загрузчиков и типов загрузки позволяет глубже понять
взаимодействие аппаратного и программного обеспечения, а также процессы,
происходящие на самых ранних этапах запуска компьютера.

\hypertarget{ux430ux43aux442ux443ux430ux43bux44cux43dux43eux441ux442ux44c}{%
\chapter{Актуальность}\label{ux430ux43aux442ux443ux430ux43bux44cux43dux43eux441ux442ux44c}}

Актуальность темы обусловлена широким распространением современных
технологий UEFI, Secure Boot, виртуализации и мультизагрузочных систем.
В условиях растущих требований к безопасности и совместимости
программного обеспечения, знание устройства загрузки ОС становится
необходимым для обеспечения стабильной и безопасной работы систем.
Особенно важно это в контексте администрирования, кибербезопасности и
восстановления после сбоев.

\hypertarget{ux43eux431ux44aux435ux43aux442-ux438-ux43fux440ux435ux434ux43cux435ux442-ux438ux441ux441ux43bux435ux434ux43eux432ux430ux43dux438ux44f}{%
\chapter{Объект и предмет
исследования}\label{ux43eux431ux44aux435ux43aux442-ux438-ux43fux440ux435ux434ux43cux435ux442-ux438ux441ux441ux43bux435ux434ux43eux432ux430ux43dux438ux44f}}

Объект исследования: процесс загрузки современных операционных систем.
Предмет исследования: виды загрузки ОС (BIOS, UEFI, MBR, GPT, загрузчики
Windows Boot Manager, GRUB, boot.efi), а также особенности и отличия в
процедурах загрузки Windows, Linux и macOS.

\hypertarget{ux43fux440ux430ux43aux442ux438ux447ux435ux441ux43aux430ux44f-ux437ux43dux430ux447ux438ux43cux43eux441ux442ux44c}{%
\chapter{Практическая
значимость}\label{ux43fux440ux430ux43aux442ux438ux447ux435ux441ux43aux430ux44f-ux437ux43dux430ux447ux438ux43cux43eux441ux442ux44c}}

Практическая значимость исследования заключается в возможности
применения полученных знаний при диагностике, установке, настройке и
восстановлении операционных систем. Изучение этапов загрузки помогает
эффективно решать проблемы несовместимости оборудования, сбоев
загрузчиков, а также корректно настраивать мультизагрузочные и сетевые
конфигурации.

\hypertarget{ux43fux440ux43eux431ux43bux435ux43cux430}{%
\chapter{Проблема}\label{ux43fux440ux43eux431ux43bux435ux43cux430}}

Многие пользователи и студенты, работающие с компьютерами, не понимают,
какие процессы происходят от момента включения питания до появления
интерфейса ОС. Отсутствие этих знаний затрудняет устранение неполадок и
настройку систем. Исследование призвано восполнить этот пробел,
предоставив систематизированные сведения о типах и механизмах загрузки
ОС.

\hypertarget{ux43fux43eux441ux442ux430ux43dux43eux432ux43aux430-ux446ux435ux43bux438-ux438-ux437ux430ux434ux430ux447}{%
\chapter{Постановка цели и
задач}\label{ux43fux43eux441ux442ux430ux43dux43eux432ux43aux430-ux446ux435ux43bux438-ux438-ux437ux430ux434ux430ux447}}

\hypertarget{ux446ux435ux43bux44c-ux440ux430ux431ux43eux442ux44b}{%
\section{Цель
работы:}\label{ux446ux435ux43bux44c-ux440ux430ux431ux43eux442ux44b}}

Провести обзор и анализ видов загрузки операционных систем, описать их
принципы, этапы и особенности реализации в различных программных средах
(Windows, Linux, macOS). \#\# Гипотеза: Анализ и сравнение различных
видов загрузки операционных систем позволит лучше понять принципы их
функционирования и повысить эффективность администрирования и поддержки
вычислительных систем. Задачи: 1. Дать определение процессу загрузки
операционной системы. 2. Рассмотреть архитектуру BIOS и UEFI, их
различия и взаимодействие с MBR/GPT. 3. Описать работу основных
загрузчиков (Windows Boot Manager, GRUB, boot.efi). 4. Изучить различные
типы загрузки --- холодную, тёплую, безопасную, многозагрузочную,
сетевую и восстановительную. 5. Сравнить особенности загрузки Windows,
Linux и macOS. 6. Подчеркнуть роль современных технологий (Secure Boot,
PXE, Recovery Mode) в повышении безопасности и устойчивости систем.

\hypertarget{ux432ux432ux435ux434ux435ux43dux438ux435}{%
\chapter{Введение}\label{ux432ux432ux435ux434ux435ux43dux438ux435}}

Процесс загрузки операционной системы --- один из ключевых этапов работы
компьютера, определяющий переход от инициализации аппаратного
обеспечения к запуску программной среды. От правильности его выполнения
зависит стабильность, безопасность и производительность всей системы.
Современные операционные системы --- Windows, Linux и macOS ---
используют различные механизмы загрузки, основанные на технологиях BIOS
и UEFI, а также на схемах разметки MBR и GPT. Эти технологии определяют,
как именно выполняется поиск и запуск загрузчика --- специальной
программы, передающей управление ядру ОС.

Изучение видов загрузки важно для понимания архитектуры современных
вычислительных систем, принципов работы загрузчиков (таких как Windows
Boot Manager, GRUB, boot.efi) и особенностей технологий, обеспечивающих
надёжность и защиту --- в частности, Secure Boot и режимов
восстановления. Тема работы актуальна, поскольку знание процессов
загрузки необходимо для администрирования, диагностики и обслуживания
компьютерных систем. Цель исследования --- рассмотреть основные виды
загрузки операционных систем и их роль в обеспечении стабильной и
безопасной работы вычислительных устройств.

\hypertarget{ux43eux441ux43dux43eux432ux43dux44bux435-ux43fux43eux43dux44fux442ux438ux44f}{%
\chapter{Основные
понятия}\label{ux43eux441ux43dux43eux432ux43dux44bux435-ux43fux43eux43dux44fux442ux438ux44f}}

Загрузка (booting) -- процесс инициализации аппаратуры компьютера и
запуска операционной системы. Он начинается с включения питания:
прошивка материнской платы (BIOS или UEFI) выполняет POST (тестирование
оборудования) и передаёт управление на загрузочную программу. В
наследуемом BIOS после POST читается MBR -- главный загрузочный сектор
диска (до 512 байт), который содержит небольшой код загрузчика и таблицу
разделов. Этот код обычно запускает загрузчик второго этапа (например,
GRUB или NTLDR). В современных системах на смену BIOS пришёл UEFI --
32/64-битная прошивка, которая работает с таблицей разделов GPT. UEFI
хранит загрузчик на специальном разделе EFI (ESP) в FAT32, а путь к нему
прописан в переменных NVRAM. Иными словами, «BIOS использует MBR, а UEFI
-- GPT». В режиме UEFI плата считывает GPT, находит раздел(ы) типа EFI и
загружает из них указанный файл-загрузчик.

В системах с BIOS (Legacy) процедура такая: прошивка находит активный
диск, читает MBR и выполняет хранимый в нём код. Этот код загружает
загрузчик первого уровня (до 446 байт), который обычно сразу переходит
на основной загрузчик (например, GRUB) на системном разделе. В режиме
UEFI прошивка из NVRAM узнаёт путь к EFI-приложению (например,
`\EFI\Microsoft\Boot\bootmgfw.efi' для Windows или
\EFI\ubuntu\grubx64.efi для Linux) и запускает его. UEFI устраняет
ограничения MBR: поддерживает диски свыше 2 ТБ и до 128 разделов. После
запуска загрузчика (Windows Boot Manager, GRUB, Apple boot.efi и т.д.)
происходит передача управления ядру ОС. Так, в Windows BIOS-или
UEFI-загрузчик активирует Boot Manager, который согласно записям BCD
находит Winload.exe (или winload.efi), загружает ядро ntoskrnl.exe и
необходимые драйверы. В Linux типовой загрузчик GRUB (Grand Unified
Bootloader) читает файл конфигурации (обычно /boot/grub/grub.cfg) и
показывает меню с установленными системами. GRUB поддерживает
многоплатформенную загрузку: он может напрямую грузить Linux-ядро и
initrd, а также «по цепочке» передавать управление другим загрузчикам
(например, загрузчику Windows). Для Mac на Intel-платах после UEFI
запускается файл boot.efi (базовый загрузчик macOS), который загружает
предварительно связанное ядро prelinkedkernel. На новых Mac с чипом
безопасности T2 Secure Enclave реализация Secure Boot соответствует
протоколам Apple, но суть схожа с обычной UEFI-загрузкой: проверяется
подпись загрузчика перед загрузкой системы. В системе после загрузчика
ОС инициализируется уже ядро операционной системы и запускаются
системные службы или пользовательские процессы. Так, Windows после
загрузки ядра запускает процесс Smss.exe, который инициализирует сессии,
настраивает подсистемы, и лишь затем разворачивается обычная работа ОС.

\hypertarget{bios-ux438-uefi-mbr-ux438-gpt}{%
\chapter{BIOS и UEFI: MBR и GPT}\label{bios-ux438-uefi-mbr-ux438-gpt}}

BIOS -- устаревшая прошивка: 16-битная, ограниченная 1 МБ оперативки во
время старта. При классическом BIOS POST завершается, когда найдена
загрузочная запись. BIOS считывает MBR, где в первых 446 байтах
находится первичный загрузчик. В MBR же хранятся записи о 4 разделах
(максимум) и метка 0xAA55 (подпись). Затем загрузчик из MBR запускает
раздел, помеченный активным. В Linux/GRUB этот код обычно загружает
«этап 1.5», умеющий читать файловую систему, а затем полноценный второй
этап (Grub2), который отображает меню и запускает ядро. Из-за
ограничений MBR поддерживаемый размер диска -- до \textasciitilde2 ТБ, а
число основных разделов -- 4. UEFI (Unified Extensible Firmware
Interface) -- современный стандарт, замена BIOS. Работает в 32/64-битном
режиме, позволяет использовать гораздо больше памяти при инициализации.
UEFI хранит загрузчики как файлы на специальном разделе EFI System
Partition (FAT32), а сам загрузчик запускается по записи в NVRAM с
указанием диска и пути. В отличие от BIOS, UEFI может непосредственно
загрузить ОС без промежуточного кода MBR. При этом используется таблица
GPT (GUID Partition Table), которая поддерживает до 9.4 зеттабайт
дискового пространства и 128 разделов.

\hypertarget{ux43eux441ux43dux43eux432ux43dux44bux435-ux437ux430ux433ux440ux443ux437ux447ux438ux43aux438-ux43eux441}{%
\chapter{Основные загрузчики
ОС}\label{ux43eux441ux43dux43eux432ux43dux44bux435-ux437ux430ux433ux440ux443ux437ux447ux438ux43aux438-ux43eux441}}

\hypertarget{windows-boot-manager}{%
\section{Windows Boot Manager}\label{windows-boot-manager}}

В Windows после POST используется специальный загрузчик Windows Boot
Manager (Bootmgr или bootmgfw.efi в UEFI). На системном разделе
находится папка \EFI\Microsoft\Boot, где лежит этот файл. При старте
BIOS/UEFI выполняет Bootmgr, который обращается к хранилищу BCD (Boot
Configuration Data) и предлагает выбрать ОС (в случае мультизагрузки)
или продолжить загрузку. Затем Bootmgr запускает Winload.exe
(NT-носитель) или winload.efi, который загружает ядро Windows и основные
драйверы. Процесс загрузки Windows описывается так: «BIOS считывает MBR
и запускает диспетчер загрузки Windows», затем «Диспетчер загрузки
Windows находит и запускает загрузчик Windows (Winload.exe)», после чего
«Загрузчик ОС загружает основные драйверы для ядра, а затем запускает
само ядро». Таким образом, последовательность Windows (Legacy BIOS) --
POST → MBR → Bootmgr → Winload → ядро. В UEFI-протоколе
последовательность аналогична, но вместо MBR прошивка напрямую запускает
EFI-приложение Bootmgr.

\hypertarget{grub-ux438-linux}{%
\section{GRUB и Linux}\label{grub-ux438-linux}}

В Linux и других UNIX-подобных ОС стандартным загрузчиком является GRUB
(GRand Unified Bootloader) версии 2. GRUB гибок: может загружать ядра
Linux, ядро FreeBSD, Solaris, а через «chainloader» -- и Windows.
Конфигурация GRUB (обычно /boot/grub/grub.cfg) задаёт пункты меню,
параметры ядра и initrd. GRUB помещается в MBR (в BIOS-системах) или в
ESP (в UEFI-системах) и выводит текстовое меню при загрузке.
Пользователь может вручную редактировать параметры ядра непосредственно
в GRUB перед запуском (например, указать runlevel или задающий init).
GRUB также поддерживает протокол UEFI, то есть в GPT-системах он
устанавливается в папке EFI как, например, \EFI\ubuntu\grubx64.efi.
Благодаря мультизагрузочной природе GRUB часто используется на системах
с несколькими ОС -- он может выбирать между ядрами разных дистрибутивов
и Windows.

\hypertarget{ux437ux430ux433ux440ux443ux437ux447ux438ux43a-macos}{%
\section{Загрузчик
macOS}\label{ux437ux430ux433ux440ux443ux437ux447ux438ux43a-macos}}

На Macintosh (Intel) нет BIOS -- используется прошивка EFI, близкая по
принципам к UEFI. При старте Mac выполняется BootROM, затем загружается
boot.efi с диска. В Apple-платформах загрузчик называется
iBoot/boot.efi. В Intel-Mac без чипа T2 прошивка UEFI просто загружает
файл boot.efi с системного раздела, который затем загружает ядро macOS
(файл prelinkedkernel). На новых Mac с чипом Apple T2 или Apple Silicon
используется цепочка доверенной загрузки: встроенный BootROM проверяет
подпись iBoot, затем UEFI проверяет boot.efi и immutablekernel,
обеспечивая полный Secure Boot для macOS. В случае проблем Mac
автоматически входит в режим восстановления (Recovery). Кроме того, при
включении Mac можно удерживать клавишу Option (Alt) для вызова
встроенного Boot Manager, который позволяет выбрать между разными
системами (например, между macOS и Windows в Boot Camp) и внешними
дисками.

\hypertarget{ux442ux438ux43fux44b-ux437ux430ux433ux440ux443ux437ux43aux438}{%
\chapter{Типы
загрузки}\label{ux442ux438ux43fux44b-ux437ux430ux433ux440ux443ux437ux43aux438}}

Загрузка ОС может различаться по контексту запуска. Основные виды
загрузки:

\begin{enumerate}
\def\labelenumi{\arabic{enumi}.}
\item
  Холодная (cold boot) -- полное включение питания. При холодной
  загрузке компьютер включается из состояния «выключен», выполняется
  полный POST и инициализация оборудования. Эта процедура длиннее, так
  как аппаратуре нужно полностью инициализироваться.
\item
  Тёплая (warm boot) -- перезагрузка без полной остановки питания,
  например при нажатии Reset или команды «Перезагрузить» в системе. При
  тёплой загрузке часть начальных процедур (например, тест POST)
  пропускается, что позволяет загрузиться быстрее. Слово «тёплый старт»
  часто используется в англоязычных описаниях как warm boot (аналог
  «перезагрузка»).
\item
  Многозагрузочная конфигурация (dual/multi-boot) -- наличие на одном ПК
  нескольких операционных систем. При такой настройке сначала
  устанавливают несколько ОС на разные разделы диска или диски, а затем
  при старте появляется меню выбора. Например, GRUB может в начале
  загрузки предложить Ubuntu или Windows, а Windows Boot Manager --
  Windows или другую версию ОС. Dual boot позволяет иметь, скажем, Linux
  и Windows на одном компьютере. Это же позволяет загружать разные
  версии ОС для тестирования. Dual Boot буквально означает
  «многозагрузка»: «это многозагрузочная конфигурация (используется для
  выбора ОС при старте)». В этом режиме при запуске появляются меню
  (Boot Manager или GRUB), где пользователь выбирает систему. Для
  корректного мультизагрузочного запуска загрузчики часто используют
  цепочку: например, GRUB «по цепочке» передаёт управление загрузчику
  Windows (NTLDR или Bootmgr).
\end{enumerate}

\hypertarget{ux431ux435ux437ux43eux43fux430ux441ux43dux430ux44f-ux437ux430ux433ux440ux443ux437ux43aux430-secure-boot}{%
\chapter{Безопасная загрузка (Secure
Boot)}\label{ux431ux435ux437ux43eux43fux430ux441ux43dux430ux44f-ux437ux430ux433ux440ux443ux437ux43aux430-secure-boot}}

Это технология UEFI, обеспечивающая проверку цифровой подписи
загружаемых модулей. Secure Boot -- это часть спецификации UEFI: при
включении ПК прошивка проверяет электронную подпись каждого
EFI-приложения (включая драйверы опций и загрузчики ОС) по ключам из
хранилища системы. Запускаются только доверенные, подписанные файлы.
Стандарт Secure Boot изначально введён в стандарты Windows 8 и далее,
чтобы предотвратить загрузку неподписанного или вредоносного ПО на
ранних этапах. Если подпись незнакомая или нарушена, загрузка
блокируется. При этом в NVRAM хранятся базы доверенных сертификатов (db,
KEK, PK) для проверки. Таким образом Secure Boot защищает цепочку
загрузки ОС: если какой-то модуль скомпрометирован, система
автоматически переходит в режим восстановления и не загружает ОС.

\hypertarget{ux441ux435ux442ux435ux432ux430ux44f-ux437ux430ux433ux440ux443ux437ux43aux430-pxe-netboot}{%
\chapter{Сетевая загрузка (PXE,
NetBoot)}\label{ux441ux435ux442ux435ux432ux430ux44f-ux437ux430ux433ux440ux443ux437ux43aux430-pxe-netboot}}

загрузка системы по сети с удалённого сервера без локальных носителей. В
среде PXE (Preboot eXecution Environment) сетевой адаптер содержит
мини-загрузчик, который через UDP/BOOTP/TFTP скачивает программу сетевой
загрузки (Network Boot Program) и передаёт ей управление. Грубо:
BIOS/UEFI направляет управление на NIC, тот по DHCP/BOOTP получает IP и
адрес TFTP, загружает начальный загрузчик (pxelinux.0 или аналог). Затем
тот, в свою очередь, может загрузить ядро Linux или Windows PE по сети.
Сетевая загрузка обычно используется для массового развёртывания ОС (по
PXE-серверу) или тонких клиентов (NetBoot на Mac). Сетевой загрузчик
действует как посредник: ПК получает инструкции и образ ОС через сеть,
что позволяет администратору централизованно управлять загрузкой. \#
Загрузка в режиме восстановления (Recovery Boot)

специальный режим для восстановления системы. Каждая ОС имеет
собственный механизм: 1. Windows Recovery Environment (WinRE). Это
мини-ОС на базе Windows PE, предустановленная в разделе восстановления.
WinRE позволяет восстанавливать систему, возвращать компьютер к рабочему
состоянию или выполнять диагностику. По умолчанию WinRE автоматически
грузится после двух неудачных стартов Windows или может быть вызван
вручную (например, через «Дополнительные параметры загрузки» или
загрузочную флешку). WinRE содержит инструменты: автоматическое
восстановление, командную строку, утилиты диагностики, сброс в исходное
состояние и т.п. 2. macOS Recovery. На Mac предусмотрен встроенный
раздел восстановления. «Режим восстановления macOS --- это встроенная
система восстановления Mac». При загрузке с удерживаемой комбинацией
клавиш (Command+R или других) запускается Recovery Assistant, где
доступны Утилиты восстановления диска, переустановки macOS,
восстановления из Time Machine и т.д. В этом режиме можно исправить
диск, переустановить систему, восстановить файлы, а также настроить
параметры безопасности. 3. Linux single-user (режим
однопользовательского восстановления). Многие дистрибутивы Linux имеют
режим восстановления (Runlevel 1, «single-user mode»). В нём система
загружается в однопользовательском режиме без запуска большинства служб
и сетевых сервисов. Этот режим удобно использовать для обслуживания
(например, ручного запуска fsck) или решения проблем. В GRUB к ядру
обычно передают параметр single или emergency, чтобы войти в режим
суперпользователя без GUI. В отличии от Windows или Mac, это не
отдельный раздел, а специальный режим работы ядра.

\hypertarget{ux437ux430ux43aux43bux44eux447ux435ux43dux438ux435}{%
\chapter{Заключение}\label{ux437ux430ux43aux43bux44eux447ux435ux43dux438ux435}}

В ходе проведённого исследования были рассмотрены ключевые этапы и
механизмы процесса загрузки операционных систем, начиная от работы
базовых систем ввода-вывода (BIOS и UEFI) и заканчивая запуском ядра и
пользовательской среды. Изучение данного процесса позволило выявить,
насколько важна организация загрузки для стабильного функционирования и
безопасности компьютерных систем. Современные технологии загрузки
развиваются в сторону повышения надёжности и защиты. Если традиционные
BIOS и MBR обеспечивали лишь базовую последовательность инициализации,
то переход к UEFI и GPT принёс гибкость, поддержку больших объёмов
данных и внедрение механизмов Secure Boot, препятствующих запуску
вредоносного кода на ранних стадиях. Это особенно важно в эпоху
активного распространения сетевых угроз и атак на уровне прошивки. 1.
Особое внимание уделено загрузчикам --- таким как GRUB, LILO, Windows
Boot Manager и boot.efi, --- которые играют роль посредников между
аппаратной частью и операционной системой. Их функциональность
обеспечивает выбор между несколькими ОС, восстановление после сбоев и
гибкую настройку параметров запуска. 2. Также были рассмотрены основные
виды загрузки --- холодная, тёплая, сетевая, безопасная и
восстановительная. Каждая из них применяется в конкретных условиях и
решает свои задачи: от повседневной работы пользователя до развёртывания
корпоративных систем и обслуживания серверных инфраструктур. 3.
Практическая значимость изучения данного вопроса заключается в том, что
понимание принципов загрузки позволяет пользователю и специалисту по
информационным технологиям эффективно устранять неполадки,
оптимизировать процесс установки и настройки операционных систем, а
также повышать уровень защиты данных. 4. Таким образом, можно сделать
вывод, что процесс загрузки --- это не просто технический этап включения
компьютера, а сложная система взаимодействий программных и аппаратных
компонентов. От правильной работы каждого звена зависит успешное
функционирование всей системы. Знание механизмов загрузки необходимо не
только инженерам и администраторам, но и каждому пользователю,
стремящемуся глубже понимать работу современных технологий.

\hypertarget{ux43bux438ux442ux435ux440ux430ux442ux443ux440ux430-ux438-ux438ux441ux442ux43eux447ux43dux438ux43aux438}{%
\chapter{Литература и
источники}\label{ux43bux438ux442ux435ux440ux430ux442ux443ux440ux430-ux438-ux438ux441ux442ux43eux447ux43dux438ux43aux438}}

\begin{enumerate}
\def\labelenumi{\arabic{enumi}.}
\tightlist
\item
  Бэкон, Дж. Операционные системы: параллельные и распределительные
  системы / пер. с англ. СПб.: Питер, 2004.
\item
  Гордеев, А.В. Операционные системы: учеб. для вузов. Изд. 2-е. СПб.:
  Питер, 2005.
\item
  Таненбаум Э., Вудхал А. Операционные системы. Разработка и реализация
  / пер. с англ. СПб.; М.; Нижний Новгород, Питер, 2007.
\item
  Таненбаум Э., Бос. Х. Современные операционные системы. Изд. 4-е. --
  СПб.: Питер, 2015.
\end{enumerate}

\hypertarget{refs}{}
\begin{CSLReferences}{0}{0}
\end{CSLReferences}
